% Fonti: artefatti M15/M15b schema 2 elencati in coda.
\section{QFT approssimata: protocolli, controlli ed esiti completi}
\label{sec:app_aqft}

Questa sezione approfondisce la Sezione~\ref{sec:qft_topologia}.
Distingue il pilota $N=15$, la compilazione strutturale $N=21$, le prove
rumorose a budget ridotto e il benchmark QPE. Sono esperimenti diversi:
la campagna principale M1--M4 resta limitata a $N=15$.

\subsection{Protocollo e controllo degli SWAP}

$k_{\mathrm{QPE}}$ limita la distanza fra qubit delle rotazioni
controllate nella QFT inversa finale; $k_{\mathrm{arith}}$ applica lo
stesso criterio agli addizionatori modulari di Beauregard.
Quest'ultimo non si applica al moltiplicatore specializzato per $N=15$.
Il rumore deriva da AGI e readout della snapshot \texttt{FakeSherbrooke},
con RZ virtuali e senza aggiunta separata di $T_1/T_2$. Il seed principale
è 42; i JSON registrano layout, schedule dei seed, batch, versioni e hash
della calibrazione. La scelta usa i batch di selezione e la valutazione
primaria usa batch holdout disgiunti.

La prima implementazione senza SWAP era errata perché invertiva soltanto
l'associazione dei bit alla misura. Poiché gli SWAP precedono la cascata,
occorre coniugare anche quest'ultima con l'inversione degli indici.
Gli output errati in \path{artifacts/_superseded_swap_bug/} non sono
evidenza corrente. Con la variante corretta, i bracci con e senza SWAP
del pilota $N=15$ hanno gli stessi conteggi di gate, profondità e
conteggi di successo: il transpiler assorbe già la permutazione nel layout.

\subsection{Pilota su \texorpdfstring{$N=15$}{N=15}}

Ogni configurazione usa 8192 shot in otto batch, quattro per selezione
e quattro per holdout. La tabella include tutti i gradi; le righe
duplicate senza SWAP sono omesse.

\begin{table}[!htbp]
\centering
\small
\caption{M15 su $N=15$: griglia completa dei gradi QPE.}
\label{tab:aqft_n15_completa}
\begin{tabular}{rrrr}
\toprule
$k_{\mathrm{QPE}}$ & ECR & profondità & $P_{\mathrm{holdout}}$ \\
\midrule
1 & 279 & 1141 & 0.6279 \\
2 & 305 & 1242 & 0.5422 \\
3 & 328 & 1330 & 0.5210 \\
4 & 360 & 1414 & 0.4771 \\
5 & 365 & 1391 & 0.4607 \\
6 & 360 & 1306 & 0.4749 \\
7 (piena) & 362 & 1351 & 0.4792 \\
\bottomrule
\end{tabular}
\end{table}

Il grado scelto sulla selezione è 1; il contrasto holdout è $+0.1487$,
IC Newcombe 95\% $[+0.1273;+0.1699]$, con 83 ECR risparmiate.
Il caso è favorevole perché $r=4$ produce fasi esatte in due bit:
non va generalizzato all'aritmetica per $N=21$.

\subsection{Struttura e prove rumorose su \texorpdfstring{$N=21$}{N=21}}

La ricognizione iniziale contiene 84 configurazioni di sola compilazione,
non 84 simulazioni rumorose. A QFT finale piena, il conteggio ECR iniziale
passa da 46\,940 per aritmetica piena a 22\,631 per
$k_{\mathrm{arith}}=1$. I circuiti ricompilati per la sensibilità
rumorosa hanno invece 49\,661 e 23\,178 ECR: conteggi ottenuti da
compilazioni differenti non vanno mescolati.

Il pilota nominale esegue soltanto il braccio troncato con 256 shot,
due batch e 128 shot holdout, ottenendo $0.2813$. Mancando il riferimento
pieno nello stesso pilota, non permette un contrasto completo.
I tre livelli successivi usano 256 shot per braccio in quattro batch,
con 128 shot holdout. La tabella confronta sempre la variante
aritmetica troncata con quella piena, anche quando la selezione
preferisce quest'ultima.

\begin{table}[!htbp]
\centering
\small
\caption{M15 su $N=21$: sensibilità al fattore di rumore.
Differenza troncata meno piena, QFT finale piena e
$k_{\mathrm{arith}}=1$ nel braccio troncato.}
\label{tab:aqft_n21_sensibilita}
\begin{tabular}{rcccl}
\toprule
fattore & piena & troncata & differenza & IC 95\% \\
\midrule
0.03  & 0.2813 & 0.2891 & $+0.0078$ & $[-0.1020;+0.1174]$ \\
0.01  & 0.2734 & 0.2734 & $0.0000$ & $[-0.1084;+0.1084]$ \\
0.003 & 0.3828 & 0.2734 & $-0.1094$ & $[-0.2205;+0.0056]$ \\
\bottomrule
\end{tabular}
\end{table}

Al fattore 0.003 la selezione sceglie la QFT piena. Per questo il
contrasto primario fra configurazione scelta e riferimento pieno
nel JSON vale zero: non è il contrasto fra troncata e piena della
tabella, ricalcolato dai conteggi holdout con lo stesso metodo Newcombe.
Tutti e tre gli intervalli includono lo zero; nessun confronto dimostra
un vantaggio della variante troncata. Il campione ridotto limita la
discriminazione di effetti piccoli. Questi test esplorativi non
estendono a $N=21$ il confronto M1--M2.

\paragraph{Il margine a miscela è una stima condizionata.}
La stima di circa tre punti percentuali usa un modello ausiliario
\begin{equation}
 P_{\mathrm{mix}}(G,e)=P_{\mathrm{floor}}+
 (P_{\mathrm{ideal}}-P_{\mathrm{floor}})(1-e)^G,
 \label{eq:aqft_miscela}
\end{equation}
con $P_{\mathrm{ideal}}=0.4667/0.4013$ per piena/troncata e
$G=49\,661/23\,178$. Assume che ogni evento d'errore conduca al
pavimento uniforme; l'AGI usata per $e$ è un proxy e non identifica la
probabilità fisica di tale evento. Il margine dipende da questa ipotesi:
non è una misura di fedeltà, una curva rumorosa simulata o un limite
superiore dimostrato per ogni canale. Rese ideali e stima a miscela
sono riportate nel README sperimentale; conteggi e prove rumorose
sono registrati nei JSON.

\subsection{Benchmark QPE: tutti i nove contrasti}

Il successo è la misura della migliore approssimazione della fase a
$n$ bit, definita da
\begin{equation}
 y_{\mathrm{target}}=\operatorname{round}(\phi 2^n)\bmod 2^n.
 \label{eq:aqft_qpe_target}
\end{equation}
Ogni grado usa 4096 shot in otto batch, quattro di selezione e quattro
holdout; il grado è scelto sulla selezione. Il confronto mantiene lo
stesso layout e la stessa schedule dei seed tra i bracci.

\begin{table}[!htbp]
\centering
\footnotesize
\setlength{\tabcolsep}{3pt}
\caption{M15b: gradi scelti e contrasti holdout completi.
IC Newcombe al 95\%; i nove punti sono esplorativi e non
ricevono una correzione simultanea.}
\label{tab:aqft_qpe_completa}
\begin{tabular}{rlrcccl}
\toprule
$n$ & fase & $k$ & scelta & piena & $\Delta$ & IC 95\% \\
\midrule
6 & $1/6$ & 2 & 0.3755 & 0.3599 & $+0.0156$ & $[-0.0139;+0.0451]$ \\
6 & $21/2^6$ & 2 & 0.5576 & 0.5059 & $+0.0518$ & $[+0.0212;+0.0822]$ \\
6 & $1/3$ & 2 & 0.3564 & 0.3418 & $+0.0146$ & $[-0.0145;+0.0438]$ \\
8 & $1/6$ & 3 & 0.2861 & 0.2598 & $+0.0264$ & $[-0.0009;+0.0536]$ \\
8 & $85/2^8$ & 3 & 0.4014 & 0.3545 & $+0.0469$ & $[+0.0172;+0.0765]$ \\
8 & $1/3$ & 3 & 0.2930 & 0.2500 & $+0.0430$ & $[+0.0157;+0.0701]$ \\
10 & $1/6$ & 2 & 0.2354 & 0.1641 & $+0.0713$ & $[+0.0469;+0.0956]$ \\
10 & $341/2^{10}$ & 2 & 0.3564 & 0.2310 & $+0.1255$ & $[+0.0977;+0.1530]$ \\
10 & $1/3$ & 2 & 0.2280 & 0.1733 & $+0.0547$ & $[+0.0302;+0.0791]$ \\
\bottomrule
\end{tabular}
\end{table}

I nove contrasti hanno stime positive, sei intervalli escludono zero.
I risparmi ECR sono 25, 28 e 78 per $n=6,8,10$.
Il grado più favorevole fra quelli provati è 2 o 3; un layout per taglia,
una snapshot e tre fasi non dimostrano generalizzazione ad altri dispositivi.
La scelta del grado resta empirica sui batch di selezione: i conteggi e
la calibrazione possono suggerire candidati, ma questa campagna non valida
un selettore che prescinda dalle esecuzioni.

\subsection{Provenienza e stato degli artefatti}

La radice è \path{Extra/experiments/M15_qft_topologia/artifacts/}.
Le fonti correnti, tutte schema 2, sono:
\begin{itemize}
 \item \path{v1_20260830/results_M15_qft_N15_v1_20260830_151350.json};
 \item \path{v1_20260830/results_M15_qft_N21_v1_20260830_152226.json}
       (sola struttura);
 \item \path{pilota/results_M15_qft_N21_v1_20260831_120126.json};
 \item \path{sensibilita/results_M15_qft_N21_v1_20260831_135104.json};
 \item \path{sensibilita/results_M15_qft_N21_v1_20260831_144129.json};
 \item \path{sensibilita/results_M15_qft_N21_v1_20260831_152704.json};
 \item \path{benchmark/results_M15b_qpe_v1_20260831_163523.json}.
\end{itemize}
Checkpoint e log documentano avanzamento e controlli accessori;
non sono campagne indipendenti. Gli output della variante SWAP errata
restano esclusi.
